%%
%% 開志専門職大学 ICT活用総合実習 最終報告書サンプル
%%
%% 情報処理学会研究報告テンプレートをベースに，
%% ゼミ報告書用のカスタムスタイル (kaishi-report.sty) を適用しています．
%%
%% コンパイル方法:
%%   platex sample-report.tex
%%   pbibtex sample-report    （参考文献を使う場合）
%%   platex sample-report.tex （相互参照の解決）
%%   platex sample-report.tex （相互参照の解決）
%%   dvipdfmx sample-report
%%
%% または latexmk を使う場合:
%%   latexmk -pdfdvi sample-report.tex
%%

%% noauthor: 1枚目下の著者名・所属を非表示にする
\documentclass[submit,techrep,noauthor]{ipsj}

%% ゼミ報告書用カスタムスタイル（ヘッダ・フッタを上書き）
\usepackage{kaishi-report}

\usepackage[dvipdfmx]{graphicx}
\usepackage{latexsym}
\usepackage{url}

%% ===== 報告書の設定 =====
%% 年度を指定（ヘッダに反映されます）
\nendo{2026}
%% 指導教員（デフォルト: 鈴木源吾．変更する場合のみ指定）
% \advisor{鈴木源吾}

\begin{document}

%% タイトル（和文のみ）
\title{○○を用いた△△の分析と考察}

%% 著者名（タイトル下に表示されます．所属の脚注は表示されません）
\author{開志 太郎}{Taro Kaishi}{}

%% 概要（和文のみ）
\begin{abstract}
本報告では，○○という課題に対して△△の手法を用いた取り組みについて述べる．
まず，背景として□□の現状と課題を整理し，
それに対する解決方法として△△を提案する．
提案手法を実装・適用した結果，◇◇の改善が確認された．
最後に，本取り組みの成果をまとめ，今後の課題について述べる．
\end{abstract}

\maketitle

%%=====================================================================
\section{はじめに}
%%=====================================================================

近年，○○の分野において△△が注目されている．
本実習では，□□という課題に着目し，
△△の手法を用いた解決を試みた．

本報告の構成は以下の通りである．
第\ref{sec:background}章では背景と課題を述べ，
第\ref{sec:method}章では課題の解決方法を説明する．
第\ref{sec:result}章では結果と考察を示し，
第\ref{sec:conclusion}章でまとめと今後の課題を述べる．

%%=====================================================================
\section{背景と課題}
\label{sec:background}
%%=====================================================================

%% ここに研究の背景と課題を記述する．
%% 先行研究や関連技術への言及も含める．

\subsection{背景}

○○の分野では，近年□□が広く利用されるようになっている\cite{okumura}．
しかし，△△に関する課題が指摘されており，
特に◇◇の点で改善の余地がある．

\subsection{課題}

本実習で取り組む具体的な課題は以下の通りである．

\begin{itemize}
\item 課題1：○○のデータが□□の形式で提供されており，そのままでは分析が困難である
\item 課題2：既存の手法では△△の精度が十分でない
\end{itemize}

%%=====================================================================
\section{課題の解決方法}
\label{sec:method}
%%=====================================================================

%% ここに提案手法・実装方法を記述する．
%% 図表を用いて分かりやすく説明する．

前章で述べた課題に対し，本実習では以下の方法で解決を試みた．

\subsection{データの前処理}

○○のデータに対して，以下の前処理を実施した．

\begin{enumerate}
\item □□形式のデータを△△形式に変換
\item 欠損値の補完（◇◇の手法を使用）
\item 正規化処理
\end{enumerate}

\subsection{分析手法}

前処理後のデータに対して，△△の手法を適用した．
具体的には，以下の手順で分析を行った．

%% 図の挿入例（実際の図ファイルに置き換えてください）
% \begin{figure}[tb]
%   \includegraphics[width=\columnwidth]{../figures/method-overview.pdf}
%   \caption{提案手法の概要}
%   \label{fig:method}
% \end{figure}

%%=====================================================================
\section{結果と考察}
\label{sec:result}
%%=====================================================================

%% ここに実験・分析の結果と，それに対する考察を記述する．
%% 表や図を活用する．

\subsection{結果}

提案手法を適用した結果を表\ref{tab:result}に示す．

\begin{table}[tb]
\caption{分析結果の比較}
\label{tab:result}
\hbox to\hsize{\hfil
\begin{tabular}{l|cc}\hline\hline
 & 既存手法 & 提案手法 \\\hline
指標1 & 0.75 & 0.82 \\
指標2 & 0.68 & 0.79 \\\hline
\end{tabular}\hfil}
\end{table}

\subsection{考察}

表\ref{tab:result}の結果から，
提案手法は既存手法と比較して指標1で約9\%，
指標2で約16\%の改善が見られた．
この改善は，△△の処理が□□に寄与したためと考えられる．

一方で，◇◇のケースでは十分な改善が得られなかった．
これは，○○の特性により△△が十分に機能しなかったためと推測される．

%%=====================================================================
\section{まとめと今後の課題}
\label{sec:conclusion}
%%=====================================================================

本報告では，○○の課題に対して△△の手法を用いた取り組みについて述べた．
提案手法により，指標1・指標2ともに改善が確認され，
□□の有効性を示すことができた．

今後の課題として，以下が挙げられる．

\begin{itemize}
\item ◇◇のケースへの対応
\item より大規模なデータでの検証
\item □□の自動化
\end{itemize}

%%=====================================================================
%% 参考文献
%%=====================================================================

\begin{thebibliography}{9}

\bibitem{okumura}
奥村晴彦：改訂第5版\LaTeXe 美文書作成入門，
技術評論社(2010).

%% 以下，参考文献を追加してください．
%% BibTeX を使う場合は thebibliography 環境の代わりに
%% \bibliographystyle{ipsjunsrt}
%% \bibliography{references}
%% を使用してください．

\end{thebibliography}

\end{document}
